\documentclass[a4paper,14pt]{report} %размер бумаги устанавливаем А4, шрифт 14 пунктов
\usepackage[T2A]{fontenc}
\usepackage[english,russian]{babel} % включаем русский язык
\usepackage[utf8x]{inputenc} % работаем с форматом  utf-8
\usepackage{tocloft}% изменяем параметры оглавления
\usepackage{graphicx} % библиотека для работы с графикой
\graphicspath{ {./img/} } % путь к папке с изображениями
\usepackage{amssymb,amsfonts,amsmath,mathtext,cite,enumerate,float} %подключаем нужные пакеты расширений
\usepackage{cite} % Красивые ссылки на литературу
\usepackage{mathtext}
\usepackage{caption} % описания
\usepackage{color} %% это для отображения цвета в коде
\usepackage{listings} %% собственно, это и есть пакет listings
\linespread{1.3}
\usepackage{geometry} % Меняем поля страницы
\geometry{left=2cm}% левое поле
\geometry{right=1.5cm}% правое поле
\geometry{top=1cm}% верхнее поле
\geometry{bottom=2cm}% нижнее поле
%%% Стили %%%
\bibliographystyle{BibTeX-Styles/utf8gost71u}    % Оформляем библиографию по ГОСТ 7.1 (ГОСТ Р 7.0.11-2011, 5.6.7)
\cftsetpnumwidth{3pt} % точки прям до цифр в оглавлении

\makeatletter
\renewcommand{\@biblabel}[1]{#1.} % Заменяем библиографию с квадратных скобок на точку:
\makeatother

\renewcommand{\theenumi}{\arabic{enumi}}% Меняем везде перечисления на цифра.цифра
\renewcommand{\labelenumi}{\arabic{enumi})}% Меняем везде перечисления на цифра.цифра
\renewcommand{\theenumii}{.\arabic{enumii}}% Меняем везде перечисления на цифра.цифра
\renewcommand{\labelenumii}{\arabic{enumi}.\arabic{enumii}.}% Меняем везде перечисления на цифра.цифра
\renewcommand{\theenumiii}{.\arabic{enumiii}}% Меняем везде перечисления на цифра.цифра
\renewcommand{\labelenumiii}{\arabic{enumi}.\arabic{enumii}.\arabic{enumiii}.}% Меняем везде перечисления на цифра.цифра
\renewcommand{\thesection}{\arabic{section}.}
\renewcommand{\thesubsection}{\arabic{section}.\arabic{subsection}.}
\newcommand{\tocsecindent}{\hspace{7mm}}
\newcommand*{\No}{\textnumero}
\definecolor{gray}{rgb}{0.5,0.5,0.5}
\DeclareCaptionFont{white}{\color{white}} %% это сделает текст заголовка белым
%% код ниже нарисует серую рамочку вокруг заголовка кода.
\DeclareCaptionFormat{listing}{\colorbox{gray}{\parbox{\textwidth}{#1#2#3}}}
\captionsetup[lstlisting]{format=listing,labelfont=white,textfont=white}
\addto\captionsrussian{\def\refname{Список литературы}}
\long\def\comment{}% многострочный комментарий

% документ
\begin{document}
\long\def\/*#1*/{}
	% Титульный лист (ГОСТ Р 7.0.11-2001, 5.1)
% титульный лист
\newpage
\begin{titlepage}
\thispagestyle{empty}

{\large
\begin{center}
%\thesisOrganization
	\textsl{Федеральное государственное бюджетное образовательное учреждение высшего профессионального образования \\* }	
	\vspace{2em}
\end{center}
}
%\vspace{0pt plus4fill} %число перед fill = кратность относительно некоторого расстояния fill, кусками которого заполнены пустые места
\IfFileExists{img/bmstu-logo.png}
{
\begin{minipage}[t]{0.13\textwidth}
  \centering\raisebox{\dimexpr \topskip-\height}{%
  \includegraphics[width=\textwidth]{bmstu-logo.png}}
  \label{fig0}
  \end{minipage}\hfill
\begin{minipage}[t]{0.65\textwidth}
\begin{center}
\large  \textsl{\textbf {\\ «Московский государственный технический университет имени Н.Э. Баумана» \\* (МГТУ им. Н.Э. Баумана) \\* }}
  \end{center}
\end{minipage}
}

\vspace{2em}

\hrulefill

\begin{center}
 	\vspace{0pt plus2fill} %число перед fill = кратность относительно некоторого расстояния fill, кусками которого заполнены пустые места
	{\large
	Факультет  <<Информатика и системы управления>>\\
	Кафедра  <<Программное обеспечение ЭВМ и информационные технологии>>
	}
\end{center}

{\Large
\begin{center}
	\textbf{Отчет по лабораторной работе №6\\* по курсу: \\* } «Функциональное и Логическое программирование» \\
\end{center}
}
%
\vspace{0pt plus4fill} %число перед fill = кратность относительно некоторого расстояния fill, кусками которого заполнены пустые места

\begin{center}
%Научный руководитель:
\newlength{\ML}
\hfillСтудент группы ИУ7-63Б:\hspace{0.2cm}\underline{М.\,В.~Заколесник}\\
\hfill\hspace{1.7cm} (Фамилия И.О.)\\
\hfillПреподователь:\hspace{0.2cm}\underline{Н.\,Б.~Толпинская}\\
\hfill\hspace{1.7cm} (Фамилия И.О.)\\
%\supervisorRegalia

%\supervisorFio
\end{center}
%

\vspace{0pt plus4fill} %число перед fill = кратность относительно некоторого расстояния fill, кусками которого заполнены пустые места
{\centeringМосква, 2019\par}
\end{titlepage} % титульный лист
	\newpage
	\tableofcontents % это оглавление, которое генерируется автоматически
	% задание и сама лабораторная работа
\newpage
\section*{Задание}
\addcontentsline{toc}{section}{\tocsecindent{Задание}}
\Large{Лабораторная работа №5}

\begin{itemize}
	\item Если запустить интерпретатор и ввести:
		\begin{itemize}
			\item (setf  a 2)
			\item (setf b 3)
			\item (setf c 5)
		\end{itemize}
	\item Что будет на экране, если вводить:
		\begin{itemize}
			\item a
			\item b
			\item c
			\item ‘a
			\item ‘(+ a c)
			\item (a)
			\item (eval ‘a)
		\end{itemize}
	\underline{Добавив:}\\
	\textbf{(defun a() ‘b) }\\
	\textbf{(defun b() 4)}
	\item Что будет на экране, если вводить:
		\begin{itemize}
			\item a
			\item b 
			\item c
			\item ‘a
			\item ‘(+ a c)
			\item (a)
			\item (eval ‘a)
			\item (a)
			\item (+ a a)
			\item (+ (b) b)
			\item (b b b)
		\end{itemize}
	\underline{Добавив:}\\
	\textbf{(setf a b)}
	\item Что будет на экране, если вводить:
		\begin{itemize}
			\item a
			\item b 
			\item c
			\item `a
			\item `(+ a c)
			\item (a)
			\item (eval `a)
		\end{itemize}
	\underline{Добавив:}\\
	\textbf{(setf c b)}
	\item Что будет на экране, если вводить:
		\begin{itemize}
			\item a
			\item b
			\item c
			\item `a
			\item `(+ a c)
			\item (a)
			\item (eval `a)
		\end{itemize}
	\underline{Добавив:}\\
	\textbf{(defun a(x y) (+ x y))}
	\item Что будет на экране, если вводить:
		\begin{itemize}
			\item a
			\item b
			\item c
			\item `a
			\item `(+ a c)
			\item (a)
			\item (eval `a)
			\item (a a a)
			\item (a b a)
		\end{itemize}
	\hrulefill 
\begin{itemize}
	\item (defun b(x y) (setf b (+ b 1)) (* x y b) ) 
\end{itemize}
	\hrulefill
	%	\hrulefill
	\item Что будет на экране, если вводить:
		\begin{itemize}
			\item a
			\item b
			\item (b b b)
		\end{itemize}
	\item Напишите функцию, которая вычисляет катет по гипотенузе и другому катету.
\end{itemize}
	% ответы на вопросы
\newpage
\section*{Ответы на вопросы}
\addcontentsline{toc}{section}{\tocsecindent{Ответы на вопросы}}
\begin{enumerate}
    \item\Large{Как представляются списки в ОП? Выписать и дать определение основным элементам языка Lisp.}\\
         Списочная ячейка состоит из двух частей, полей first и rest. Каждое из полей содержит указатель. Указатель может ссылаться на другую списочную ячейку или на некоторый другой Lisp объект, как, например, атом. Указатели между ячейками образуют как бы цепочку, по которой можно из предыдущей ячейки попасть в следующую и так, наконец, до атомарных объектов. Каждый известный системе атом записан в определённом месте памяти лишь один раз.\\
    \underline{Основные элементы языка:}
        \begin{itemize}
            \item Атомы:
                \begin{itemize}
                    \item Символ
                    \item Символьная константа
                    \item Число(целые, действительные, рациональные)
                    \item Логические(T, Nil)
                \end{itemize}
            \item Точечные пары:
                \begin{itemize}
                    \item (<атом> . <атом>)
                    \item (<точ. пара> . <атом>)
                    \item (<атом> . <точ. пара>)
                    \item (<точ. пара> . <точ. пара>)
                \end{itemize}
            \item S-выражения:
                \begin{itemize}
                    \item <атом> | <точ.пара>
                \end{itemize}
            \item Списки:
                \begin{itemize}
                    \item (<s-выражение> . <список>)
                    \item (<пустой список>) $\equiv$ Nil
                \end{itemize}
        \end{itemize}
    \item\Large{Как выполняются CAR и CDR?}
        \begin{itemize}
            \item car - переходит по car указателю и возвращает голову
            \item cdr - переходит по cdr указателю и возвращает хвост(остаток)
        \end{itemize}
    \item Классификация функций:
        \begin{enumerate}
            \item Конструкторы(cons, list)
                \begin{itemize}
                    \item \textbf{cons}- двухъаргументная функция, создаёт одну списковую ячейку и расставляет указатели;
                        \\ \textbf{(cons 'a 'b) -> (a.b)} 
                        \\ \textbf{(cons 'a '(b)) -> (a b)}
                    \item \textbf{list}- создаёт столько списковых ячеек, сколько аргументов
                        \\ \textbf{(list 'a 'b) -> (a b)}
                    \\
                    cons работает быстрее, чем list, но может быть организован не список
                \end{itemize}
            \item Селекторы(car, cdr)
                \begin{itemize}
                    \item \textbf{car}- переходит по car указателю и возвращает голову
                    \item \textbf{cdr}- переходит по cdr указателю и возвращает хвост(остаток)
                \end{itemize}
            \item Предикаты(позволяющие определить структуру элементов)
                \begin{itemize}
                    \item \textbf{atom}- возвращает T, если аргумент- атом, иначе Nil;
                    \item \textbf{listp}-  возвращает T, если аргумент- список, иначе Nil;
                    \item \textbf{consp}- представлена ли пара в виде списковых ячеек;
                    \item \textbf{numberp}- возвращает T, если аргумент- число, иначе Nil;
                    \item \textbf{symbolp}- возвращает T, если аргумент- не число, иначе Nil;
                \end{itemize}
            \item Сравнение(позволяющие сравнивать элементы):
                \begin{itemize}
                    \item \textbf{базовая eq}- сравнивает два символьных атома(указатели на них) на их тождественность возвращает Nil или T, применима только для символов.
                        \\ \textbf{(eq T T) -> T}
                    \item \textbf{eql}- сравнивает два числа(по форме их представления)
                        \\ \textbf{(eql 3 3) -> T}
                        \\ \textbf{(eql 3 3.0) -> Nil}
                    \item \textbf{=}- сравнивает два числа
                        \\ \textbf{(= 3 3) -> T}
                        \\ \textbf{(= 3 3.0) -> T}
                    \item \textbf{equal}- сравнивает как и eql, но также может сравнивать списки
                    \item \textbf{equalp}- сравнивает все варианты представления, но долго работает
                \end{itemize}
            \underline{Для чисел:}
                \begin{itemize}
                    \item \textbf{oddp}- проверка на нечетность
                    \item \textbf{evenp}- проверка на четность
                \end{itemize}
        \end{enumerate}
\end{enumerate}

% ответ
%\\
	% список литературы
\newpage
\addcontentsline{toc}{section}{\tocsecindent{Список литературы}}

\begin{thebibliography}{}
    \bibitem{litlink1}  Толпинская Н.Б.  -  Курс лекций по "Функциональному и Логическому программированию"[Текст], Москва 2019 год.
    \bibitem{litlink2}  Городняя Л.В.  -  Основы функционального программирования. Курс лекций : учеб. пособие для вузов / Городняя Л. В. - М. : Интернет-Университет Информационных Технологий, 2004. - 272 с. - (Основы информационных технологий). - Библиогр.: с. 269-272 и в конце кн. - ISBN 5-9556-0008-6.
    \bibitem{link3} Пол Грэм - ANSI Common Lisp, 2012. - 448 с. - Символ-Плюс и в конце кн. ISBN 0-13-370875-6.
\end{thebibliography}

\end{document}